\phantomsection
\chapter*{MasterChef Gagal}
\addcontentsline{toc}{section}{MasterChef Gagal --- Azzahra Ramadhani}
\penulis{Azzahra Ramadhani}


\awalcerpen{R}{encana Ayah buat} gaya-gayaan masak nasi goreng spesial hari Minggu berubah jadi bencana nasional di dapur kami. Dengan percaya diri tingkat tinggi dan celemek bunga-bunga punya Ibu, Ayah melempar bumbu halus ke wajan panas ala koki profesional di TV. Masalahnya, bukannya bau harum gurih yang tercium, asap hitam pekat malah membumbung tinggi memenuhi plafon rumah. Ibu yang lagi santai nonton sinetron langsung lari ke dapur sambil bawa kain basah, siap-siap mau memadamkan kebakaran yang ternyata cuma bersumber dari bumbu gosong ciptaan suami tercintanya.


Kepanikan makin tidak terkendali waktu alarm asap tetangga sebelah ikut-ikutan bunyi gara-gara asapnya keluar lewat jendela. Adik yang panik bukannya ngambil air, malah sibuk mengipasi asapnya pakai kipas angin mini, yang bikin asap hitam itu makin merata ke seluruh ruangan. Sementara itu, Kakak dengan tenangnya cuma nyangkut di pintu dapur sambil sibuk merekam muka Ayah yang sudah cemong kena noda gosong, siap diunggah ke grup keluarga besar buat bahan ejekan seminggu ke depan. Ayah cuma bisa tersenyum pasrah sambil memegang sodet, berusaha membela diri kalau warna hitam di wajan itu adalah "efek karamelisasi tingkat tinggi".


Ujung-ujungnya, nasi goreng eksperimen Ayah yang bentuknya sudah mirip arang itu berakhir di tempat sampah. Dapur yang tadi rapi berubah wujud jadi kayak bekas tempat latihan perang, lengkap dengan bekas kerak hitam di mana-mana. Alih-alih makan enak hasil masakan Ayah, kami sekeluarga akhirnya duduk berjejer di teras depan rumah sambil makan mie instan cup yang dibeli Kakak dari warung madura sebelah. Ayah yang kena omel Ibu sepanjang makan cuma bisa nyengir sambil janji tidak akan pernah menyentuh kompor lagi seumur hidupnya.


Pas lagi asyik Ngobrol, Bibi adiknya Ibu yang rumahnya cuma beda dua gang tiba-tiba datang mau mengembalikan cetakan kue. Dia langsung berhenti di pagar gara-gara mencium bau sangit yang tajam banget. Begitu melihat muka Ayah yang masih belang-belang hitam dan dapur rumah yang masih mengepul, Bibi cuma bisa menggeleng-geleng sambil ketawa tipis. Ibu yang masih cemberut langsung menceritakan kronologi kejadiannya dari awal, sementara Ayah cuma bisa diam sambil terus menyeruput kuah mie instannya.


Nggak lama setelah Bibi pulang, HP Ayah mulai bergetar terus-terusan. Ternyata video rekaan Kakak sudah menyebar di grup keluarga besar dan panen reaksi. Semua sibuk membalas pakai stiker tertawa, sedangkan Nenek langsung mengirim pesan suara panjang yang isinya menasihati Ayah supaya nggak usah gaya-gayaan masak lagi kalau cuma bikin repot orang rumah. Ayah yang makin tersudut cuma bisa menghela napas pasrah, tahu kalau kejadian malam ini bakal terus dibahas tiap ada acara kumpul keluarga.


Selesai makan, kami semua langsung masuk buat membersihkan dapur yang berantakan. Ibu membagi tugas dengan adil Adik bertugas membuka semua jendela lebar-lebar biar udaranya berganti, Kakak menyapu sisa-sisa bumbu yang berceceran di lantai, dan Ayah dapat tugas paling berat yaitu menggosok kerak hitam di wajan sampai bersih. Walau bau gosongnya masih tersisa sedikit sampai besok harinya, setidaknya suasana rumah kembali tenang, meski kami yakin janji Ayah buat nggak menyentuh kompor lagi biasanya cuma bertahan sebentar.

\jedahias
